تصدرت التحديات المرتبطة \textbf{بارتفاع تكاليف الطاقة والمياه} قائمة المعوقات، إذ يؤدي الارتفاع المستمر في أسعار الطاقة والمياه إلى زيادة تكاليف الإنتاج خاصة لأنشطة كثيفة الاستهلاك لهما، مما يمثل عبئا إضافيا على الشركات. تلتها التحديات المرتبطة بارتفاع \textbf{معدلات التضخم} رغم تراجعها بشكل ملحوظ بالمقارنة بالربع السابق. وجاءت في المرتبة الثالثة التحديات المتعلقة \textbf{بإجراءات التعامل مع الجهات الحكومية} والتي تمثل مصدر معاناة لمجتمع الأعمال نتيجة بطء الإجراءات، والروتين، وتعامل بعض الموظفين، فضلا عن تعدد موظفي الضبطية القضائية في معظم الجهات الحكومية، بما يفتح مجال لانتشار الفساد والمصروفات غير الرسمية. أما في المرتبة الرابعة، فجاءت تحديات \textbf{صعوبة الحصول على العمالة المؤهلة} وهو ما يشكل عبئا إضافيا على الشركات من خلال زيادة إنفاقها على التدريب والتأهيل، نتيجة الفجوة ما بين مخرجات منظومة التعليم والتدريب واحتياجات سوق العمل.
